\chapter{Fazit und Handlungsempfehlung}

\section{Zusammenfassung der Ergebnisse}
Um die Kontrolle über interne Daten zu behalten und gleichzeitig den steigenden Compliance-Anforderungen gerecht zu werden, ist eine technische Pseudonymisierung bei der Nutzung externer \ac{LLM}-\acp{API} für Unternehmen heute unerlässlich \footcite{staab2025sokprivacyparadox}. Mit einem hybriden Ansatz, der die deterministische Präzision von \ac{Regex} für strukturierte Daten mit der semantischen Flexibilität von \ac{NER} für unstrukturierte Kontexte kombiniert, können die Vorteile beider Ansätze genutzt werden \footcite{naseer2025recap}.Trotzdem sollte beachtet werden,
dass eine vollständige Pseudonymisierung oft schwierig ist und die
Qualität der Ergebnisse beeinträchtigen kann. Deshalb sind kontinuierliche Optimierungen notwendig, um ein zu striktes Maskierungsverfahren zu vermeiden \footcite{wang2025revisitingtradeoffs}.

\section{Handlungsempfehlung für die Praxis}
Um ein solches Modell in der Praxis umzusetzen, gibt es verschiedene Anbieter, welche Open-Source-Lösungen anbieten. Zu nennen sind hierbei insbesondere Microsoft Presidio, Skyflow oder Privacera \footcite{vignati2024towardsmiddleware}. Dabei kann zum Beispiel ein Stufenmodell eingesetzt werden, welches eine von Grund auf strikte Maskierung sukzessiv je nach Anwendungsfall lockert. Wichtig dabei sind vor allem regelmäßige Audits der Maskierungs-Logs, um Fehlklassifizierungen der Modelle zu finden \footcite{zhang2025surveyprivacypreservation}.